\section{Introduction}
\label{sec:rs-intro}

The deployment of autonomous multi-agent systems into consequential environments — agricultural resource management, financial orchestration, clinical decision support, and critical infrastructure — has exposed a structural gap that policy alone cannot close. Conventional architectures can describe what agents do; they cannot architecturally guarantee that every agent's authority traces to a human principal. As agent populations spawn recursively, modify state, and execute decisions with real-world consequences, the question of who authorized each action, through what chain of delegation, and with what scope constraints has become the central unsolved problem in agentic infrastructure.

We introduce the \textbf{Recursive Spawning Protocol} (RSP): a structural mechanism built on the \bxthree{} three-layer accountability architecture that makes the delegation chain from every active agent to a human principal a runtime artifact — immutable, cryptographically enforced, and architecturally verifiable — rather than a forensic reconstruction attempted after the fact. RSP is the second named pillar of the \bxthree{} Framework, complementing Loop Isolation, Spatial Firewall, Sandbox Gate, and the Bailout Protocol. Together, these five pillars constitute the upstream accountability guarantee: that \bxthree{} systems fail upward into human consciousness, never downward into autonomous chaos.

\subsection{The Autonomous Drift Problem}

The failure mode RSP targets is \emph{autonomous drift}: the condition in which an agent continues to operate with no traceable authorization path to a human principal. We model a multi-agent system as a rooted directed acyclic graph (DAG) of agent nodes. Each node $n$ carries a \emph{parent pointer} $\alpha(n)$ referencing the node that authorized its provisioning. The root of any authorized chain is a human principal $h \in \mathcal{H}$. A \emph{spawn chain} $C = \langle n_0, n_1, \ldots, n_k \rangle$ satisfies $\alpha(n_i) = n_{i-1}$ for all $1 \leq i \leq k$ and $n_0 = h$.

\medskip
\noindent\textbf{Definition 1 (Orphaned Agent).} An agent $n$ is \emph{orphaned} if its parent node $\alpha(n)$ has either (a) \emph{terminated} — exited operational state through normal completion, abnormal failure, or administrative shutdown — or (b) \emph{become unreachable} — is non-responsive to any defined communication protocol and this condition persists beyond the heartbeat timeout. Orphaning severs the delegation chain at the parent node. The agent persists in operational state but is disconnected from its chain of authorization.

\medskip
\noindent\textbf{Definition 2 (Drifted Agent).} An agent $n$ is \emph{drifted} if the chain $C(n) = \langle n_0, n_1, \ldots, n \rangle$ does not originate at a human principal. Equivalently: there exists no human anchor $h \in \mathcal{H}$ such that the chain terminates at $n_0 = h$. A drifted agent may be orphaned — its parent has disappeared — or may have a living parent whose own delegation chain terminates at a non-human origin. Drift is a property of the chain's origin, not of intermediate connectivity.

\medskip
\noindent\textbf{Definition 3 (Drift Triad).} The \emph{drift triad} is the conjunction of three independent conditions that produces the RSP's target failure state:
\begin{equation}
\text{DriftTriad}(n) = \langle \text{orphaned}(n),\ \text{drifted}(n),\ \text{operating}(n) \rangle.
\label{eq:intro-drift-triad}
\end{equation}
An agent enters the drift triad when all three conditions hold simultaneously: it is structurally disconnected from its parent, its delegation chain does not reach any human anchor, and it continues to execute its event loop. The drift triad is the failure condition the Recursive Spawning Protocol is designed to make architecturally impossible.

Drift is \emph{heritable}: once an agent enters the drifted state, all descendants inherit the condition absent an explicit structural intervention at spawn time. A single undetected drift event at depth $d$ with branching factor $b$ produces up to $b^d$ drifted agents. This exponential amplification — a single unchecked drift event at depth $d=10$ with $b=2$ yields $2^{10}=1024$ drifted agents before any human supervisor is aware of the initial failure — is what makes drift, rather than any single-agent failure, the central threat in recursive spawning systems.

\subsection{Why Existing Approaches Are Insufficient}

The conventional architectural response to unbounded agent spawning is the \emph{depth limit}: the system enforces a maximum chain depth $D_{\max}$, rejecting any spawn that would exceed this bound. A parallel conventional approach is the \emph{time-to-live} (TTL) constraint, which caps an agent's lifetime. Both are necessary resource-governance tools; neither addresses the root cause of drift.

Consider a chain $C = \langle n_0, n_1, \ldots, n_k \rangle$ where $n_0$ is a human principal and $k < D_{\max}$. The depth constraint is satisfied by construction. Now apply scope creep: each node $n_i$ for $i \geq 1$ successively expands its operating scope $R(n_i)$ beyond $R(n_{i-1})$ without Purpose Layer authorization. After $k$ successive expansions, $R(n_k)$ may be entirely disjoint from $R(n_0)$. Chain depth is $k \leq D_{\max}$. The chain has a human anchor at $n_0$. Yet $n_k$ is drifted by definition: its operating scope is not a descendant refinement of its human anchor's intent. A depth limit constrains the \emph{number} of spawn generations; it says nothing about the \emph{scope} of each generation's authorized actions.

The drift prevention requirement is formally a conjunction of two constraints:
\begin{equation}
\text{Drift Prevention} \iff \bigl(\text{depth}(C) \leq D_{\max}\bigr)\ \land\ \bigl(\forall i:\ R(n_{i+1}) \subseteq R(n_i)\bigr).
\label{eq:intro-drift-prevention}
\end{equation}
Depth limits enforce only the first conjunct. A system that enforces $D_{\max}$ but not the scope refinement constraint is immune to unbounded spawn cascades but remains fully vulnerable to scope creep drift. TTL constrains an agent's \emph{lifetime}, not its \emph{scope within that lifetime}. An agent with a 60-second TTL and unrestricted scope expansion authority can produce consequences exceeding its human anchor's intent within seconds of spawning.

Logging and audit infrastructure can partially reconstruct the chain after the fact, but retroactive reconstruction is categorically different from runtime enforcement. In consequential environments where unauthorized actions may be irreversible — a water rights diversion executed, a financial instruction dispatched, a medical record modified — runtime enforcement is a structural requirement, not a preference. The deeper reason depth limits and TTL fail is that they address the wrong dimension: they constrain how many spawn generations may exist or how long a node may operate; they say nothing about what each generation is authorized to do. Without scope refinement as a structural invariant — enforced not by policy convention but by protocol constraint — chain length is bounded while scope divergence is unbounded.

\subsection{Core Insight: Immutable Parent Pointer Chain}

The Recursive Spawning Protocol's core insight is that autonomous drift is structurally impossible when the parent pointer $\alpha(n)$ is immutable — established once at provisioning and unalterable by any actor under any circumstances.

An immutable parent pointer eliminates the conditions that produce drift by construction. Chain opacity cannot arise: the child's parent pointer always traces to the node that authorized its provisioning, regardless of subsequent events. The delegation chain is not a mutable reference — it is a first-class runtime artifact established by the Fact Layer write protocol and hash-chained into the tamper-evident Forensic Ledger. Scope mutation cannot produce drift without Purpose Layer authorization: the Fact Layer pre-commit hook rejects any spawn event where the proposed child scope is not a descendant refinement of the parent's scope. The subset relationship $R(n_{i+1}) \subseteq R(n_i)$ is not a policy convention — it is a structural invariant enforced at every spawn event.

The immutable parent pointer is not a storage policy or an access control configuration. It is a write protocol constraint enforced by the Fact Layer at the protocol level — before any write is executed, not after it is observed. Any attempt to overwrite $\alpha(n)$ after provisioning is rejected regardless of the requestor's identity, privileges, or justification. There is no administrative override; there is no privileged actor capable of modifying the chain after the fact. The immutability is structural, not contractual.

This enables runtime verification that no conventional architecture can provide: at any point during execution, the system can construct the complete delegation chain for any agent $n$, verify that the chain terminates at a human anchor $h(n_0)$, and confirm that each node's operating scope is a descendant refinement of the human anchor's intent — without depending on the agent's self-reporting, the completeness of external logging infrastructure, or the honesty of any machine actor in the chain.

RSP adds three constraints to conventional agentic runtimes: an immutable parent pointer established at provisioning and protected by the Fact Layer write protocol; a Purpose Layer authorization check at every spawn event requiring scope refinement justification and spawn necessity declaration; and a Bounds Engine depth bound enforced as a structural gate. Together, these constraints make autonomous drift architecturally impossible — not statistically unlikely, not conventionally prevented, but structurally impossible as a matter of protocol enforcement.

\subsection{Paper Roadmap}

The remainder of this paper is organized as follows. Section~\ref{sec:rs-background} situates RSP within the research traditions it draws on — agent lifecycle management, accountability and provenance in distributed systems, fixed versus mutable delegation structures, and hierarchical task decomposition — and establishes the \bxthree{} Framework as the minimal sufficient substrate for RSP's accountability guarantees. Section~\ref{sec:rs-drift} provides the formal characterization of autonomous drift, defines the drift triad and its constituent failure sub-modes, catalogs the four concrete failure modes that arise from the three enabling structural conditions, and proves formally that depth limits alone are insufficient to prevent drift. Section~\ref{sec:rs-protocol} specifies the Recursive Spawning Protocol in full: its Purpose Layer validation gate, its Bounds Engine depth management, its Fact Layer immutable pointer and forensic ledger, and its chain verification and convergence criteria. Section~\ref{sec:rs-integration} maps each RSP component to its governing \bxthree{} layer and demonstrates the structural necessity of the three-layer mapping. Section~\ref{sec:rs-correctness} formally states and proves the drift prevention, chain integrity, and termination correctness properties. Section~\ref{sec:rs-limitations} examines the structural costs of RSP's guarantees and identifies directions for future work. Section~\ref{sec:rs-conclusion} concludes with a summary of RSP's contribution to the accountability of multi-agent AI systems.
